% HSK 电工杯 LaTeX 通用起稿模板
% Template lineage/history is not the current Skill release; current release is declared by core/bootstrap.yaml and release carriers.
% 编译：python scripts/render_paper.py final_latex --profile diangong --clean
% 当届规则、页数、封面和指定章节要求始终优先于本通用模板。
\documentclass[12pt,a4paper]{ctexart}

\usepackage[margin=2.54cm]{geometry}
\usepackage{amsmath,amssymb,amsthm,bm}
\usepackage{graphicx}
\usepackage{float}
\usepackage{booktabs}
\usepackage{longtable}
\usepackage{multirow}
\usepackage{listings}
\usepackage{xcolor}
\usepackage{hyperref}
\usepackage{caption}
\usepackage{subcaption}
\usepackage{enumitem}
\usepackage{titlesec}
\usepackage{fancyhdr}
\usepackage{lastpage}
\usepackage[backend=biber,style=numeric,sorting=none]{biblatex}
\addbibresource{references.bib}

\newcommand{\ProblemCode}{A}
\newcommand{\TeamNumber}{XXXX}
\newcommand{\PaperTitle}{电工杯论文标题}

\hypersetup{colorlinks=false,pdfborder={0 0 0}}
\lstset{
    basicstyle=\small\ttfamily,
    breaklines=true,
    frame=single,
    numbers=left,
    numberstyle=\tiny,
    showstringspaces=false
}

\pagestyle{fancy}
\fancyhf{}
\fancyhead[L]{队伍编号 \TeamNumber}
\fancyhead[R]{第 \thepage\ 页 / 共 \pageref{LastPage} 页}
\renewcommand{\headrulewidth}{0.4pt}
\titleformat*{\section}{\Large\bfseries}
\titleformat*{\subsection}{\large\bfseries}
\titleformat*{\subsubsection}{\normalsize\bfseries}

\begin{document}

\begin{titlepage}
\begin{center}
\vspace*{2cm}
{\Huge\bfseries \PaperTitle}\\[2cm]
{\Large 题号：\textbf{\ProblemCode}}\\[1cm]
{\Large 队伍编号：\textbf{\TeamNumber}}\\[3cm]
\end{center}
\end{titlepage}

% 以下仅提供通用语义骨架；正文内容、章节数量与标题按当届题面和当前项目事实调整。
\section{问题重述}
\subsection{问题背景}
% 在此概括研究背景与对象。
\subsection{问题提出}
% 在此逐问重述任务、输入、输出和约束。

\section{问题分析}
% 在此说明各问之间的数据、参数、模型或结果依赖。

\section{问题一模型建立及求解}
\subsection{模型建立}
% 在此给出变量、假设、目标函数、约束或核心关系。
\subsection{模型求解}
% 在此说明求解器、算法流程与关键前提。
\subsection{求解结果}
% 在此报告当前主结果及必要的数值证据。
\subsection{结果的分析与验证}
% 在此解释结果、验证适用性并说明结论边界。

\section{模型的评价与推广}
% 在此基于实际证据评价模型优缺点、适用范围与可推广方向。

% AI 工具使用声明是条件合规位置，固定放在参考文献之前，但默认不生成可见声明。
% 只有 config/competition_profiles.yaml 中当前电工杯当届 ai_disclosure 规则已核验、
% 且参赛队已确认当前项目真实 AI 使用事实并判定披露适用时，才取消下面示意注释并据实填写。
% 不得把模板示意原样提交，也不得默认声称参赛队使用了 AI。
% \section*{AI工具使用声明}
% <仅填写已核验规则要求/允许披露且参赛队确认属实的工具、用途、范围与责任说明。>

\printbibliography[title={参考文献}]

\appendix
\section{主要算法与复现说明}
\section{补充推导与图表}

\end{document}